\section{Erd\H{o}s Problem \#884 --- continuation}

\section{ROUND-2 OBJECTIVE}

I pursue path \textbf{(C)}: isolate and prove a \emph{fatal obstruction} to the main proof strategy from the previous round.

Round-(round\_no-1) proved a positive theorem for integers of the form
\[
n=r^a m,\qquad (r,m)=1,
\]
with constants depending on
\[
M:=\tau(m),
\]
by using only the fact that the divisor exponents lie in a union of \(M\) translates of \(\mathbb Z\). The exact remaining gap was whether one could hope to remove this \(M\)-dependence by pushing that same translate-decomposition method harder.

The new result of this round shows that this is impossible in a very strong sense: even in a highly rigid repeated-block model, any inequality of the form
\[
S_{\mathrm{all}}\le C_M\,S_{\mathrm{gap}}+B_M
\]
must satisfy
\[
C_M\gg \log M.
\]
In particular, the \(H_{M-1}\)-type loss appearing in the previous-round theorem is \emph{qualitatively sharp} in the ambient ``union of \(M\) geometric progressions'' model, even if one allows an arbitrary additive error \(B_M\).

This rules out an entire class of possible gap-closing strategies.

\section{Round-(round\_no-1) FOUNDATION USED}

I use the following Round-(round\_no-1) results as fixed input.

\begin{enumerate}
\item The notation
\[
S_{\mathrm{all}}(n):=\sum_{1\le i<j\le t}\frac{1}{d_j-d_i},
\qquad
S_{\mathrm{gap}}(n):=\sum_{1\le i<t}\frac{1}{d_{i+1}-d_i}.
\]
\item The previous-round theorem for numbers \(n=r^a m\) with \((r,m)=1\), which gives an upper bound with constants depending on \(M=\tau(m)\), obtained solely from the fact that the divisor exponents lie in a union of \(M\) translates of \(\mathbb Z\).
\item The previously checked contextual status that the original Erd\H{o}s problem remains unresolved unconditionally, while Tao's 2025 note provides a conditional disproof under a qualitative prime-tuples conjecture.
\end{enumerate}

I do \emph{not} re-prove any of those results here.

\section{NEW INSIGHT / TOOL (ROUND-2)}

The new tool is a repeated-block obstruction theorem.

The previous round already showed that arbitrary unions of \(M\) geometric progressions can force a logarithmic loss. What is new here is that one can prove the same logarithmic obstruction in a much more rigid model that mirrors the exact structure exploited by the previous proof: one takes \(M\) translates, truncates each to the same length \(L\), and arranges them in disjoint ``blocks''.

More precisely, for integers \(M\ge 2\) and \(L\ge 1\), define
\[
r_L:=e^{1/L},
\qquad
\lambda_{L,M}:=e^{1/(2ML)},
\]
and consider the increasing sequence
\[
X_{L,M}:=\{x_{k,j}: 0\le k<L,\ 0\le j<M\},
\qquad
x_{k,j}:=r_L^k\lambda_{L,M}^j=e^{(k+j/(2M))/L}.
\]
This sequence lies in a union of exactly \(M\) geometric progressions with common ratio \(r_L\), and each block
\[
\{x_{k,0},\dots,x_{k,M-1}\}
\]
is itself a short geometric progression with ratio \(\lambda_{L,M}\).

The main new theorem proves that for this model
\[
\sum_{i<j}\frac{1}{x_j-x_i}
\gg L^2M^2\log M,
\qquad
\sum_i\frac{1}{x_{i+1}-x_i}
\ll L^2M^2,
\]
so the ratio is \(\gg \log M\). Since this holds for every \(L\), any additive term \(B_M\) is asymptotically irrelevant.

This is strictly stronger than the previous-round obstruction: it shows not just that no absolute constant can work in the ambient model, but that the \emph{coefficient of \(S_{\mathrm{gap}}\)} itself must be at least of order \(\log M\).

\section{ATTACK PLAN (ROUND-2)}

The exact gap left after Round-(round\_no-1) was this:

\begin{quote}
Can one hope to remove the \(M\)-dependence from the translate-decomposition method, perhaps by refining the bookkeeping, while still using only the geometry of a union of \(M\) translates of \(\mathbb Z\)?
\end{quote}

I prove the answer is \emph{no}.

The precise claims to prove are:

\begin{enumerate}
\item For the explicit repeated-block model \(X_{L,M}\), the full reciprocal-gap energy is bounded below by \(cL^2M^2H_{\lfloor M/2\rfloor}\).
\item For the same model, the consecutive-gap reciprocal sum is bounded above by \(CL^2M^2\).
\item Consequently,
\[
\frac{\sum_{i<j}1/(x_j-x_i)}{\sum_i 1/(x_{i+1}-x_i)}\gg H_{\lfloor M/2\rfloor}\asymp \log M.
\]
\item Therefore any theorem in the ambient model of the form
\[
S_{\mathrm{all}}\le C_M S_{\mathrm{gap}}+B_M
\]
valid for all such sequences must satisfy \(C_M\gg \log M\). In particular, the previous-round \(H_{M-1}\)-type loss is qualitatively sharp.
\end{enumerate}

This overcomes the previous ambiguity: it does not merely point to a heuristic obstruction, but proves a rigorous lower bound that rules out an entire family of improvement attempts.

\section{WORK (ROUND-2)}

\subsection*{1. The repeated-block model}

Fix integers \(M\ge 2\) and \(L\ge 1\). Define
\[
r:=r_L=e^{1/L},
\qquad
\lambda:=\lambda_{L,M}=e^{1/(2ML)}.
\]
Consider the finite set
\[
X_{L,M}:=\{x_{k,j}: 0\le k<L,\ 0\le j<M\},
\qquad
x_{k,j}:=r^k\lambda^j=e^{(k+j/(2M))/L}.
\]

For each fixed \(j\), the set
\[
\{x_{k,j}:0\le k<L\}
\]
is a geometric progression with common ratio \(r\). Hence \(X_{L,M}\) is contained in a union of exactly \(M\) geometric progressions with common ratio \(r\).

Also, for each fixed \(k\), the block
\[
B_k:=\{x_{k,0},x_{k,1},\dots,x_{k,M-1}\}
\]
is a geometric progression with common ratio \(\lambda\).

Because
\[
x_{k,M-1}=e^{(k+(M-1)/(2M))/L}<e^{(k+1)/L}=x_{k+1,0},
\]
the increasing order on \(X_{L,M}\) is exactly the lexicographic order on \((k,j)\): first all points of \(B_0\), then all points of \(B_1\), and so on.

Define
\[
A_{L,M}:=\sum_{x<y\in X_{L,M}}\frac{1}{y-x},
\qquad
G_{L,M}:=\sum_{\text{consecutive }x<y\text{ in }X_{L,M}}\frac{1}{y-x}.
\]

\subsection*{2. Upper bound for the consecutive-gap sum}

\paragraph{Lemma 1.}
For all integers \(M\ge 2\) and \(L\ge 1\),
\[
G_{L,M}\le 5L^2M^2.
\]

\emph{Proof.}
There are two kinds of consecutive gaps.

\medskip
\noindent\textbf{Within-block gaps.}
For \(0\le k<L\) and \(0\le j<M-1\),
\[
x_{k,j+1}-x_{k,j}=x_{k,j}(\lambda-1)=e^{(k+j/(2M))/L}(\lambda-1).
\]
Therefore the total within-block contribution is
\[
G^{\mathrm{in}}_{L,M}
=
\sum_{k=0}^{L-1}\sum_{j=0}^{M-2}
\frac{1}{e^{(k+j/(2M))/L}(\lambda-1)}.
\]
Since
\[
\lambda-1=e^{1/(2ML)}-1\ge \frac{1}{2ML}
\]
by the elementary inequality \(e^x-1\ge x\), we obtain
\[
G^{\mathrm{in}}_{L,M}
\le 2ML\sum_{k=0}^{L-1}e^{-k/L}\sum_{j=0}^{M-2}e^{-j/(2ML)}.
\]
Now
\[
\sum_{j=0}^{M-2}e^{-j/(2ML)}\le M,
\]
and because \(e^{-1/L}\in (0,1)\),
\[
\sum_{k=0}^{L-1}e^{-k/L}
\le \frac{1}{1-e^{-1/L}}.
\]
For \(0<x\le 1\), one has \(1-e^{-x}\ge x/2\), so with \(x=1/L\),
\[
1-e^{-1/L}\ge \frac{1}{2L},
\qquad\text{hence}
\qquad
\sum_{k=0}^{L-1}e^{-k/L}\le 2L.
\]
Thus
\[
G^{\mathrm{in}}_{L,M}\le 2ML\cdot 2L\cdot M=4L^2M^2.
\]

\medskip
\noindent\textbf{Between-block gaps.}
For \(0\le k<L-1\), the gap between the last point of \(B_k\) and the first point of \(B_{k+1}\) is
\[
x_{k+1,0}-x_{k,M-1}
=e^{k/L}\left(e^{1/L}-e^{(M-1)/(2ML)}\right).
\]
Now
\[
\frac{1}{L}-\frac{M-1}{2ML}=\frac{M+1}{2ML}\ge \frac{1}{2L},
\]
so
\[
e^{1/L}-e^{(M-1)/(2ML)}
=e^{(M-1)/(2ML)}\left(e^{(M+1)/(2ML)}-1\right)
\ge \frac{M+1}{2ML}
\ge \frac{1}{2L},
\]
again using \(e^x-1\ge x\). Therefore
\[
\frac{1}{x_{k+1,0}-x_{k,M-1}}
\le 2L\,e^{-k/L}.
\]
Summing over \(k=0,\dots,L-2\),
\[
G^{\mathrm{btw}}_{L,M}
\le 2L\sum_{k=0}^{L-2}e^{-k/L}
\le 2L\cdot 2L=4L^2.
\]

Combining the two estimates,
\[
G_{L,M}=G^{\mathrm{in}}_{L,M}+G^{\mathrm{btw}}_{L,M}
\le 4L^2M^2+4L^2
\le 5L^2M^2,
\]
because \(M\ge 2\). This proves the lemma.
\hfill\(\Box\)

\subsection*{3. Lower bound for the full reciprocal-gap energy}

\paragraph{Lemma 2.}
For all integers \(M\ge 2\) and \(L\ge 1\),
\[
A_{L,M}
\ge
\frac{e^{-3/2}}{2}
L^2M^2H_{\lfloor M/2\rfloor},
\]
where
\[
H_n:=\sum_{u=1}^{n}\frac1u.
\]

\emph{Proof.}
It suffices to keep only the contributions from pairs lying in the same block. For fixed \(k\), define
\[
A_k:=\sum_{0\le i<j<M}\frac{1}{x_{k,j}-x_{k,i}}.
\]
Then
\[
A_{L,M}\ge \sum_{k=0}^{L-1}A_k.
\]
Now
\[
x_{k,j}-x_{k,i}=e^{(k+i/(2M))/L}\bigl(e^{(j-i)/(2ML)}-1\bigr),
\]
so if we write \(u=j-i\), then
\[
A_k
=
\sum_{u=1}^{M-1}\sum_{i=0}^{M-1-u}
\frac{1}{e^{(k+i/(2M))/L}\bigl(e^{u/(2ML)}-1\bigr)}.
\]
Hence
\[
A_k
=e^{-k/L}
\sum_{u=1}^{M-1}\sum_{i=0}^{M-1-u}
\frac{e^{-i/(2ML)}}{e^{u/(2ML)}-1}.
\]
Restrict to \(1\le u\le \lfloor M/2\rfloor\). For such \(u\), we have \(M-u\ge M/2\), and for all admissible \(i\),
\[
0\le \frac{i}{2ML}\le \frac{M-1}{2ML}\le \frac12,
\qquad\text{so}\qquad
e^{-i/(2ML)}\ge e^{-1/2}.
\]
Also,
\[
0<\frac{u}{2ML}\le \frac14,
\]
so the elementary bound \(e^x-1\le 2x\) for \(0\le x\le 1\) gives
\[
e^{u/(2ML)}-1\le \frac{u}{ML}.
\]
Therefore
\[
\frac{e^{-i/(2ML)}}{e^{u/(2ML)}-1}
\ge e^{-1/2}\frac{ML}{u}.
\]
Summing over \(i=0,\dots,M-1-u\) gives
\[
A_k
\ge e^{-k/L}e^{-1/2}
\sum_{u=1}^{\lfloor M/2\rfloor}(M-u)\frac{ML}{u}.
\]
Since \(M-u\ge M/2\) on this range,
\[
A_k
\ge e^{-k/L}e^{-1/2}\frac{M^2L}{2}
\sum_{u=1}^{\lfloor M/2\rfloor}\frac1u
=
 e^{-k/L}e^{-1/2}\frac{M^2L}{2}H_{\lfloor M/2\rfloor}.
\]
Now \(0\le k\le L-1\) implies \(e^{-k/L}\ge e^{-1}\), hence
\[
\sum_{k=0}^{L-1}e^{-k/L}\ge Le^{-1}.
\]
Therefore
\[
A_{L,M}
\ge
\sum_{k=0}^{L-1}A_k
\ge
Le^{-1}\cdot e^{-1/2}\frac{M^2L}{2}H_{\lfloor M/2\rfloor}
=
\frac{e^{-3/2}}{2}L^2M^2H_{\lfloor M/2\rfloor}.
\]
This proves the lemma.
\hfill\(\Box\)

\subsection*{4. Sharpness theorem for the \(M\)-progressions method}

\paragraph{Theorem 3.}
For every pair of integers \(M\ge 2\) and \(L\ge 1\), the sequence \(X_{L,M}\) defined above satisfies
\[
A_{L,M}
\ge
\frac{e^{-3/2}}{10}
H_{\lfloor M/2\rfloor}
G_{L,M}.
\]
In particular,
\[
A_{L,M}\gg (\log M)\,G_{L,M}
\qquad (M\to\infty),
\]
uniformly in \(L\).

\emph{Proof.}
This follows immediately by combining Lemma 1 and Lemma 2:
\[
A_{L,M}
\ge
\frac{e^{-3/2}}{2}L^2M^2H_{\lfloor M/2\rfloor},
\qquad
G_{L,M}\le 5L^2M^2.
\]
Hence
\[
A_{L,M}
\ge
\frac{e^{-3/2}}{10}H_{\lfloor M/2\rfloor}G_{L,M}.
\]
Since \(H_{\lfloor M/2\rfloor}\asymp \log M\), the final statement follows.
\hfill\(\Box\)

\subsection*{5. Consequence for the previous-round theorem}

\paragraph{Corollary 4.}
Suppose that for each integer \(M\ge 2\) there exist real numbers \(C_M\ge 0\) and \(B_M\ge 0\) such that every finite increasing sequence contained in a union of \(M\) geometric progressions with a common ratio satisfies
\[
\sum_{i<j}\frac{1}{x_j-x_i}
\le C_M\sum_i\frac{1}{x_{i+1}-x_i}+B_M.
\]
Then necessarily
\[
C_M\ge \frac{e^{-3/2}}{10}H_{\lfloor M/2\rfloor}
\gg \log M.
\]
In particular, no bound with \(C_M=o(\log M)\) is possible in this ambient model, regardless of the size of the additive term \(B_M\).

\emph{Proof.}
Apply the assumed inequality to the explicit sequence \(X_{L,M}\). Then
\[
A_{L,M}\le C_M G_{L,M}+B_M.
\]
By Theorem 3,
\[
\frac{e^{-3/2}}{10}H_{\lfloor M/2\rfloor}G_{L,M}
\le A_{L,M}
\le C_M G_{L,M}+B_M.
\]
Rearranging,
\[
\left(\frac{e^{-3/2}}{10}H_{\lfloor M/2\rfloor}-C_M\right)G_{L,M}
\le B_M.
\]
Now the within-block gaps alone give a matching lower growth for \(G_{L,M}\). Indeed,
\[
G_{L,M}\ge \sum_{k=0}^{L-1}\sum_{j=0}^{M-2}\frac{1}{x_{k,j+1}-x_{k,j}}
=\sum_{k=0}^{L-1}\sum_{j=0}^{M-2}\frac{1}{e^{(k+j/(2M))/L}(\lambda-1)}.
\]
Since \((k+j/(2M))/L\le 3/2\), one has \(e^{-(k+j/(2M))/L}\ge e^{-3/2}\), and since
\[
\lambda-1=e^{1/(2ML)}-1\le \frac{1}{ML}
\]
by the inequality \(e^x-1\le 2x\) for \(0\le x\le 1\), we conclude that
\[
G_{L,M}\ge e^{-3/2}L(M-1)ML.
\]
In particular \(G_{L,M}\to\infty\) as \(L\to\infty\) for fixed \(M\). Therefore the coefficient of \(G_{L,M}\) on the left must be non-positive, which gives
\[
C_M\ge \frac{e^{-3/2}}{10}H_{\lfloor M/2\rfloor}.
\]
This is exactly the claim.
\hfill\(\Box\)

\subsection*{6. Interpretation for Erd\H{o}s \#884}

The previous round proved a theorem for divisor sets of the form \(n=r^a m\) with constants depending on \(M=\tau(m)\), using only the fact that the divisor exponents lie in a union of \(M\) translates of \(\mathbb Z\). Corollary 4 shows that this type of argument is already qualitatively optimal in its ambient model:

\begin{quote}
Even allowing an arbitrary additive error term, any theorem obtained solely from the ``union of \(M\) geometric progressions'' structure must pay a coefficient at least of order \(\log M\) in front of the consecutive-gap sum.
\end{quote}

Thus the remaining difficulty in Erd\H{o}s \#884 is not a minor bookkeeping issue in the previous proof. Any eventual unconditional proof with an \emph{absolute} constant would have to use genuinely new arithmetic input about actual divisor sets, beyond the translate-decomposition geometry.

\section{ADVERSARIAL VERIFICATION}

I now check the new work against the main failure modes.

\paragraph{1. Is the ordering of \(X_{L,M}\) really lexicographic?}
Yes. The largest element of block \(B_k\) is
\[
e^{(k+(M-1)/(2M))/L},
\]
which is strictly smaller than the smallest element of the next block,
\[
e^{(k+1)/L},
\]
because \((M-1)/(2M)<1\).

\paragraph{2. Are the upper and lower bounds uniform in \(L\) and \(M\)?}
Yes. Every estimate used only the elementary inequalities
\[
e^x-1\ge x \quad (x\ge 0),
\qquad
e^x-1\le 2x \quad (0\le x\le 1),
\]
and the trivial bounds
\[
\sum_{j=0}^{M-2}e^{-j/(2ML)}\le M,
\qquad
\sum_{k=0}^{L-1}e^{-k/L}\le 2L,
\qquad
\sum_{k=0}^{L-1}e^{-k/L}\ge Le^{-1}.
\]
No hidden asymptotic assumptions were used.

\paragraph{3. Does the additive term really become irrelevant?}
Yes. For fixed \(M\), Lemma 1 gives
\[
G_{L,M}\ll L^2M^2,
\]
while the within-block gaps alone yield the lower bound
\[
G_{L,M}\gg L^2M^2.
\]
Also Lemma 2 gives
\[
A_{L,M}\gg L^2M^2H_{\lfloor M/2\rfloor}.
\]
Thus, for fixed \(M\), both \(A_{L,M}\) and \(G_{L,M}\) grow quadratically in \(L\), while an additive term \(B_M\) is independent of \(L\). Sending \(L\to\infty\) forces the claimed lower bound on \(C_M\).

\paragraph{4. Does this disprove the original Erd\H{o}s problem?}
No. The sequences \(X_{L,M}\) are model sequences in a union of \(M\) geometric progressions; they are not, in general, divisor sequences of integers. The theorem does \emph{not} settle Erd\H{o}s \#884 itself. It proves instead that the entire proof strategy based only on the translate-decomposition geometry cannot yield an absolute constant.

\paragraph{5. Is this genuinely stronger than the previous round?}
Yes. The previous round showed only that no \emph{absolute} constant can hold uniformly in the ambient model. The new theorem shows more: even after allowing an arbitrary additive error, the coefficient of \(S_{\mathrm{gap}}\) itself must be at least of order \(\log M\). This proves that the previous-round \(H_{M-1}\)-type loss is qualitatively sharp.

\section{FINAL}

\textbf{UNRESOLVED (BUT STRICTLY ADVANCED).}

The original Erd\H{o}s problem remains unresolved unconditionally. However, this round proves a sharper obstruction theorem than before:

\begin{quote}
In the ambient class of finite increasing sequences contained in a union of \(M\) geometric progressions, even bounds of the form
\[
S_{\mathrm{all}}\le C_M S_{\mathrm{gap}}+B_M
\]
require
\[
C_M\gg \log M.
\]
Thus the \(H_{M-1}\)-type coefficient from the previous round is qualitatively optimal for that entire method, and no refinement based only on the translate-decomposition geometry can produce an absolute constant.
\end{quote}

This is strict progress: it rules out a substantially narrower class of proof strategies than the previous round did, and it isolates the next barrier sharply. Any eventual unconditional proof of Erd\H{o}s \#884 must use arithmetic information about actual divisor sets that is invisible in the ambient geometric-progressions model.

\section{COMPLETION ESTIMATE}

\noindent \textbf{COMPLETION: 92\%.}

\section{REFERENCES}

\begin{enumerate}
\item T. Tao, \emph{On the sum of reciprocals of gaps between divisors}, September 9, 2025.
\item T. F. Bloom, \emph{Erd\H{o}s Problem \#884}.
\end{enumerate}
